\section{Glossario dei Termini di Dominio}

La presente sezione definisce in modo formale e non ambiguo i termini specialistici, gli attori e i concetti fondamentali impiegati nella specifica del sistema \textbf{MyAma}.

\begin{center}
\renewcommand{\arraystretch}{1.3}
\begin{longtable}{|p{4.5cm}|p{11cm}|}
\hline
\rowcolor{tableheader} \textbf{Termine} & \textbf{Definizione e Significato nel Sistema} \\ \hline
\endfirsthead

\hline
\rowcolor{tableheader} \textbf{Termine} & \textbf{Definizione e Significato nel Sistema} \\ \hline
\endhead

\textbf{MyAma} & Piattaforma digitale integrata dedicata alla gestione, pianificazione e prenotazione dei servizi di smaltimento dei rifiuti ingombranti e speciali per la città di Roma. \\ \hline

\textbf{Cittadino (Non Registrato / Visitatore)} & Utente esterno non autenticato che accede al sistema per consultare informazioni pubbliche sui servizi, tariffe, regolamenti e sedi AMA attive, e che può effettuare la registrazione. \\ \hline

\textbf{Cittadino (Registrato)} & Utente esterno autenticato che usufruisce dei servizi di MyAma per richiedere ritiri a domicilio o prenotare conferimenti in sede, monitorare lo stato delle proprie richieste, rilasciare valutazioni e ricevere notifiche. \\ \hline

\textbf{Autista AMA} & Dipendente operativo AMA incaricato dello svolgimento dei ritiri a domicilio. Consulta i ritiri assegnati, raggiunge l'indirizzo indicato dal Cittadino, effettua il carico e registra l'esito del servizio. \\ \hline

\textbf{Operatore di sede AMA} & Dipendente operativo AMA che presidia un centro di raccolta (sede/isola ecologica). Gestisce l'accoglienza del Cittadino, verifica la prenotazione e la conformità del rifiuto e registra l'esito del conferimento. \\ \hline

\textbf{Lavoratore AMA} & Categoria generale del personale operativo che comprende Autisti e Operatori di sede. Ogni lavoratore è caratterizzato da mansioni, turni di servizio e disponibilità oraria. \\ \hline

\textbf{Amministratore di sede AMA} & Figura responsabile della gestione logistica e operativa di una specifica sede AMA. Gestisce l'anagrafica del personale della sede, le disponibilità di lavoratori e mezzi, le fasce orarie e l'associazione tra sede e CAP/zone servite. \\ \hline

\textbf{Amministratore Generale AMA} & Figura direttiva a livello aziendale. Gestisce gli account degli Amministratori di sede (creazione, abilitazione e revoca) e accede a report e statistiche aggregate sui servizi erogati. \\ \hline

\textbf{Ritiro a domicilio} & Servizio logistico in cui AMA preleva il rifiuto ingombrante direttamente presso l'indirizzo indicato dal Cittadino. Richiede la verifica di copertura del CAP, l'assegnazione di un Autista e di un Veicolo compatibile con il carico. \\ \hline

\textbf{Conferimento in sede} & Servizio in cui il Cittadino si reca personalmente presso una sede AMA (centro di raccolta) per consegnare il rifiuto ingombrante, previa prenotazione di una fascia oraria. \\ \hline

\textbf{Prenotazione} & Entità informativa che formalizza la richiesta di smaltimento (a domicilio o in sede). Include dettagli sul rifiuto, indirizzo/sede, data e fascia oraria, dati del Cittadino e stato di avanzamento. \\ \hline

\textbf{Stato della prenotazione} & Condizione in cui si trova una prenotazione nel suo ciclo di vita: \textit{In attesa, Confermata, In corso, Completata, Annullata, Non eseguita}. \\ \hline

\textbf{Codice Prenotazione / Pass di Conferimento} & Identificativo univoco (codice alfanumerico o QR code) generato dal sistema alla conferma della prenotazione, utilizzato per il riconoscimento e la validazione del servizio al varco o al domicilio. \\ \hline

\textbf{Tariffa / Preventivo} & Quota economica calcolata dal sistema per l'erogazione del servizio di smaltimento, determinata in base alla tipologia, al volume del rifiuto o alla presenza di servizi speciali oltre la franchigia comunale gratuita. \\ \hline

\textbf{Itinerario di Ritiro} & Sequenza ordinata degli appuntamenti di ritiro a domicilio assegnati a uno specifico Autista e Mezzo per un dato turno lavorativo all'interno della zona di competenza. \\ \hline

\textbf{Valutazione / Feedback} & Giudizio qualitativo (espresso tramite scala di punteggio ed eventuali note descrittive) che il Cittadino può rilasciare al termine di un servizio completato per misurare gradimento e puntualità. \\ \hline

\textbf{Rifiuto ingombrante} & Bene durevole, oggetto voluminoso o materiale speciale (es. mobili, grandi elettrodomestici, RAEE) che non può essere conferito nei normali cassonetti stradali. \\ \hline

\textbf{Tipologia di rifiuto} & Classificazione merceologica e normativa del rifiuto (es. legno, metallo, RAEE, ingombranti misti), necessaria per verificare la conformità di scarico e il mezzo idoneo. \\ \hline

\textbf{Sede AMA / Centro di Raccolta} & Struttura fisica territoriale (isola ecologica) abilitata alla ricezione e allo stoccaggio temporaneo di determinate tipologie di rifiuti conferiti dai cittadini. \\ \hline

\textbf{CAP / Zona servita} & Suddivisione territoriale che delimita la competenza operativa di una sede AMA e determina se un indirizzo è coperto dal servizio di ritiro a domicilio. \\ \hline

\textbf{Veicolo / Mezzo AMA} & Automezzo aziendale impiegato per i ritiri a domicilio, caratterizzato da limiti di portata utile (peso massimo) e volume di carico. \\ \hline

\textbf{Disponibilità / Slot Orario} & Fascia temporale definita (data e intervallo orario) in cui un servizio può essere prenotato ed erogato, calcolata in base alla capienza residua dei mezzi e ai turni del personale. \\ \hline

\textbf{Assegnazione} & Associazione formale e logistica tra una prenotazione di ritiro a domicilio e le risorse operative aziendali (Autista, Veicolo, data e percorso). \\ \hline

\textbf{Esito del servizio} & Registrazione conclusiva dell'intervento da parte del personale AMA (\textit{Completato con successo, Cittadino assente, Rifiuto non conforme/respinto}). \\ \hline

\textbf{Notifica} & Comunicazione automatica generata dal sistema (email, SMS o notifica in-app) per aggiornare il Cittadino o il personale su conferme, promemoria o variazioni di stato delle prenotazioni. \\ \hline

\textbf{Report / Statistiche} & Informazioni aggregate e indicatori di performance (volumi gestiti, ritiri completati, trend per zona) consultabili dalla direzione aziendale per finalità analitiche e decisionali. \\ \hline

\textbf{Sistema} & L'applicazione software integrata MyAma nel suo complesso, comprensiva di tutti i moduli web e mobili per i diversi profili utente. \\ \hline

\textbf{Codice di Invito} & Codice alfanumerico univoco generato dagli amministratori, utilizzato per autorizzare e abilitare la registrazione del personale aziendale assegnando automaticamente il ruolo corretto. \\ \hline

\textbf{RBAC (Role-Based Access Control)} & Modello di controllo degli accessi adottato dalla piattaforma che regola i permessi di visualizzazione e interazione degli utenti in base al ruolo aziendale o utente ricoperto. \\ \hline

\textbf{Competenza Territoriale} & Regole logiche che vincolano i servizi di ritiro a domicilio e di conferimento in sede a specifiche aree geografiche (identificate tramite CAP) gestite da una determinata Sede AMA. \\ \hline

\textbf{Capacità del Veicolo} & Parametro indicante il carico logistico massimo (in peso o volume) supportato da un automezzo, utilizzato dal sistema per limitare e calcolare quanti ritiri assegnare in un singolo turno. \\ \hline

\textbf{Informativa Privacy} & Documentazione legale che esplicita le modalità di trattamento e conservazione dei dati personali, accettata obbligatoriamente da tutti gli utenti durante la procedura di registrazione. \\ \hline

\textbf{Responsive Design} & Paradigma di sviluppo dell'interfaccia utente che garantisce l'adattabilità automatica di layout e contenuti alle dimensioni dello schermo del dispositivo utilizzato (smartphone, tablet, PC). \\ \hline

\textbf{Hashing} & Metodologia crittografica unidirezionale adottata dal sistema per la memorizzazione sicura delle credenziali di accesso (password) all'interno del database. \\ \hline
\end{longtable}
\end{center}
